\documentclass[12pt]{article}

\usepackage{xcolor}
\renewcommand\thefootnote{\textcolor{red}{\arabic{footnote}}}
\usepackage[margin=1in]{geometry}
\usepackage{fancyhdr,graphicx, parskip}
\usepackage[spanish]{babel}
\usepackage{amsmath,amssymb,amsfonts,latexsym,cancel}\usepackage{rawfonts}
\usepackage{pictexwd}


\renewcommand{\headrulewidth}{0pt}
\fancyhead[L]{\includegraphics[width=3cm]{logo_uv.eps}}




\begin{document}
	
	\title{\textbf{Soluciones Tarea Obligatoria II\\ M\'etodos Matem\'aticos F\'isica II }}
	\author{\textit{$^{1}$Vicente Herrera, $^{2}$Marcelo \'Ordenes, $^{3}$Felipe Ortiz, $^{4}$Fabi\'an Trigo.}\\
		\small{\textit{$^{1,2,3,4}$Estudiantes de 3er. a\~{n}o de la Licenciatura en F\'isica}}\\
		\small {\textit{Universidad de Valpara\'iso, Facultad de Ciencias, CHILE.}}
	}
	\date{\small{\today}}
	
	\maketitle
	\thispagestyle{fancy}
\rule{112mm}{0.1mm}\\
\textbf{Problema III: Integraciones varias}\\
\rule{112mm}{0.1mm}
\begin{enumerate}
	\item [A)]
	\item [B)]
	\item [C)] " Evaluar la siguiente integral de indice 1: $J = \int_{0}^{\infty} exp(-\alpha x)cos(\beta x) dx $" para experimentar un poco y probar los brackets, podemos representar el $cos(\beta x) = Re[exp(ix\beta)]$
	
	Luego la integral J pasaria a ser Re[J']:
	$ J' = \int_{0}^{\infty} exp(-\alpha x)exp(ix\beta) dx $
	
	$ J' = \int_{0}^{\infty} exp((-\alpha + i\beta )x) dx $
	
	$z = -\alpha + i \beta$
	
	Expansion a la exponencial: $exp(zx) = \sum_{n=0}^{\infty} \Phi_n (-z)^n (x)^n $
	
	Aplicando la integral en x y convirtiendolo en un bracket:
	
	$J' = \sum_{n=0}^{\infty} \Phi_n (-z)^n < n+1 > $
	
	$J' = (\alpha - i\beta)^n \Gamma(-n) |_{n=-1}$
	
	Multiplicamos por 1, en forma del conjugado
	
	$J' = \frac{1}{\alpha - i\beta} \frac{\alpha + i\beta}{\alpha + i\beta}$
	
	$J' = \frac{\alpha + i\beta}{\alpha^2 + \beta^2}$
	
	$J = Re[J'] = \frac{\alpha}{\alpha^2 + \beta^2}$
	
\end{enumerate}



\end{document}